\exercisegroup

\begin{enumerate}
\def\labelenumi{\arabic{enumi})}
\item
  设 X 是道路连通空间，x 是 X 中的一点，\(\mathcal{U} = (\mathrm{U}_{i})\) 是 X 的开覆盖；对每个 i，令 \(c_{i}\) 为 X 中一条道路的同伦类，该道路起点为一点 \(u_i\)，且该点属于 \(U_i\)，终点为 \(x\)。记 \(G_{\mathcal{U}}\) 为 \(\pi_1(X,x)\) 的最小正规子群，使得 \(\text{Int}(c_i)(G_{\mathcal{U}})\) 包含 \(U_i\) 中以 \(u_i\) 为基点的每个圈的同伦类。则 \(G_{\mathcal{U}}\) 对容许拓扑是 \(\pi_1(X,x)\) 的开子群。
\item
  设 X 是拓扑空间，\(a\) 是 X 中的一点。若 \(\pi_1(\mathrm{X}, a)\) 的子群 H 满足：对所有 \(b \in \mathrm{X}\) 及每个道路类 \(\gamma \in \varpi_{a,b}(\mathrm{X})\)，都存在 \(b\) 的一个邻域 V，使得对每个 \(c \in \Omega_b(\mathrm{V})\)，圈 \(\gamma[c]\gamma^{-1}\) 的同伦类属于 H，则称 H 为适当子群。
\end{enumerate}

\begin{enumerate}
\def\labelenumi{\alph{enumi})}
\item
  证明 \(\pi_{1}(X,a)\) 上存在唯一的群拓扑，使适当子群恰好是开子群。
\item
  证明此拓扑比紧收敛拓扑更细，但这两个拓扑具有相同的开正规子群。
\item
  若 X 局部道路连通，则适当拓扑和紧收敛拓扑具有相同的开子群。
\end{enumerate}

\begin{enumerate}
\def\labelenumi{\arabic{enumi})}
\setcounter{enumi}{2}
\tightlist
\item
  令 C 为直径为 \(((0,0),(0,1))\) 的 \(R^{2}\) 中的圆，令 P 为所有 \(2^{n}C\)（\(n \in Z\)）的并。令 E 为将空间 \(P \times R\) 按满足 \((x,t) \sim (2x,t+1)\) 对所有 \((x,t) \in P \times R\) 都成立的最粗等价关系取商所得的空间。令 p 表示从 \(P \times R\) 到 E 的典范满射。令 a 为 P 中的点 \((0,0)\)，并置 \(e = p(a,0)\)。
\end{enumerate}

\begin{enumerate}
\def\labelenumi{\alph{enumi})}
\item
  证明 p 使 \(P \times R\) 成为群 Z 的伽罗瓦覆叠。证明 \(p_{*}(\pi_{1}(P \times \mathbf{R}, (a,0)))\) 是 \(\pi_{1}(E, e)\) 的开正规子群，且商群为 Z。
\item
  令 \(\varphi\colon\mathrm{P}\to\mathrm{P}\) 为如下定义的映射：若 \(\varphi(x)=x\) 当 \(x\in\mathrm{C}\) 时，否则 \(\varphi(x)=(0,0)\)。证明它是连续的。
\item
  证明同态 \(\varphi_{*} = \pi_{1}(\varphi, (a,0))\) 的核是 \(\pi_{1}(\mathrm{P} \times \mathbf{R}, (a,0))\) 的一个子群，并且对于紧收敛拓扑是开子群。
\item
  令 \(u: I \to C\) 为由 \(t \mapsto \frac{1}{2}(1 - \cos(2\pi t), \sin(2\pi t))\) 给出的圈，令 \(c \in \Omega_e(E)\) 为由 \(t \mapsto (u(t), 0)\) 给出的圈。证明 \(\varphi_*\) 的像是 \(\pi_1(E, e)\) 的开子群，但不包含圈 c 的同伦类。
\item
  令 \(c_{n}: I \to E\) 为由 \(c_{n}(t) = p(2^{-n}u(t), 0)\) 定义的圈，令 \(d: I \to E\) 为由 \(t \mapsto p(a, t)\) 给出的圈。证明 \(d^{n}c_{n}\overline{d}^{n}\) 严格同伦于 c。推出同伦类 \([c]\) 包含于 \(\pi_{1}(E, e)\) 的每个容许子群中。
\item
  证明 \(\pi_{1}(\mathrm{E},e)\) 上的容许拓扑严格粗于紧收敛拓扑。
\end{enumerate}

\begin{enumerate}
\def\labelenumi{\arabic{enumi})}
\setcounter{enumi}{3}
\tightlist
\item
  令 L 为 Alexandroff 半直线（TG，IV，第 49 页，习题 12），令 \(L^{*}\) 为其 Alexandroff 紧化；令 \(\omega\) 为 \(L^{*} - L\) 中的唯一点。令 A 是至少含有两个元素的集合，a 是 A 中的一点；置 \(X = L^{*} \times A\)，\(Y = X - \{(\omega, a)\}\)。
\end{enumerate}

\begin{enumerate}
\def\labelenumi{\alph{enumi})}
\item
  证明从 Y 到 X 的典范嵌入 j 是平展的、分离的，并具有道路提升性质。
\item
  证明映射 j 不能使 Y 成为 X 的覆叠。
\end{enumerate}

\begin{enumerate}
\def\labelenumi{\arabic{enumi})}
\setcounter{enumi}{4}
\tightlist
\item
  令 \(X_{0}\) 为数值平面 \(R^{2}\) 的单位圆；令 \(X_{1}\) 为连接点 \((1,0)\) 与点 \((3,1/n)\)（\(n \in N^{*}\)）的线段之并；令 \(X_{2}\) 为平面上点 \((x,y)\) 所成的集合，其中满足 \((x-2)^{2} + y^{2} = 1\) 且 \(y \leqslant 0\)；置 \(X = X_{0} \cup X_{1} \cup X_{2}\)。
\end{enumerate}

\begin{enumerate}
\def\labelenumi{\alph{enumi})}
\item
  令 \(Y = X \cup (-X)\)。令 \(p: Y \to X\) 为如下定义的映射：有 \(p(\cos(t), \sin(t)) = (\cos(2t), \sin(2t))\)（\(t \in R\)），且有 \(p(z) = p(-z) = z\)（\(z \in X_{1} \cup X_{2}\)）。证明 p 使 Y 成为 X 的二重覆叠。
\item
  令 \(Z_{2}\) 为平面上满足 \((x-1)^{2}+y^{2}=4\) 且 \(y\leqslant0\) 的点集；置 \(Z=X_{0}\cup X_{1}\cup(-X_{1})\cup Z_{2}\cup(-Z_{2})\)。令 \(q:Z\to X\) 为如下定义的映射：有 \(q(\cos(t),\sin(t))=(\cos(2t),\sin(2t))\)（\(t\in R\)）；有 \(q(z)=q(-z)=z\)（\(z\in X_{1}\)）；有 \(q(z)=q(-z)=1+\frac{1}{2}z\)（\(z\in Z_{2}\)）。证明 q 使 Z 成为 X 的二重覆叠。
\item
  证明 \(q_{*}\pi_{1}(\mathrm{Z},(1,0)) = p_{*}\pi_{1}(\mathrm{Y},(1,0))\)，但不存在连续映射 \(f\colon Z \to Y\) 使得 \(p \circ f = q\)。\(^{(4)}\)
\end{enumerate}

\begin{enumerate}
\def\labelenumi{\arabic{enumi})}
\setcounter{enumi}{5}
\tightlist
\item
  对每个整数 \(n \geqslant 1\)，置 \(x_{n} = (1/n,0)\)，并将圆记为 \(\mathrm{C}_{n}\)，其直径为 \((0,x_{n})\)。令 P 为所有空间 \(\mathrm{C}_{n}\)（\(n \in \mathbf{N}\)）的并（“夏威夷耳环”，参见 I，第 146 页，习题 5）；对每个整数 \(n \geqslant 1\)，令 \(\mathrm{P}_{n}\) 为所有空间 \(\mathrm{C}_{m}\)（\(1 \leqslant m \leqslant n\)）的并，令 \(\mathrm{P}^{n}\) 为所有空间 \(\mathrm{C}_{m}\)（\(m > n\)）的并。
\end{enumerate}

\begin{enumerate}
\def\labelenumi{\alph{enumi})}
\item
  令 n 是满足 \(\geqslant 1\) 的整数，令 \(r_{n}\) 为从 P 到 \(P_{n}\) 的映射：将 P 中的点 \(x \in P\) 在 \(x \in P_{n}\) 时映到其自身，否则映到 0。证明 \(r_{n}\) 是从 \(P_{n}\) 到 P 的典范嵌入的连续收缩映射。
\item
  证明 \(P_{n}\) 是局部可缩的。推出 P 中除原点外的每一点都有一个在该点处的可缩邻域。
\item
  证明 \(P_{n}\) 和 \(P^{n}\) 到 P 的典范嵌入诱导出从自由积 \(\pi_{1}(\mathrm{P}_{n},0)*\pi_{1}(\mathrm{P}^{n},0)\) 到 \(\pi_{1}(\mathrm{P},0)\) 的群同构。（证明 \(P^{n}\) 到补集 \(U_{n}\)（即集合 \(\{x_{1},\ldots,x_{n}\}\) 的补集）中的典范嵌入是同伦等价；然后对 P 的覆叠 \((\mathrm{P}_{n},\mathrm{U}_{n})\) 应用 IV，第 422 页的命题 2。）
\end{enumerate}

\(d)\) 令 \(\gamma \in \pi_1(\mathrm{P},0)\) 为满足 \((r_n)_*(\gamma) = \varepsilon_0\)（对每个整数 \(n \geqslant 1\)）的圈类；证明 \(\gamma = \varepsilon_0\)。

\begin{enumerate}
\def\labelenumi{\alph{enumi})}
\setcounter{enumi}{4}
\item
  从 \(\pi_1(\mathrm{P},0)\) 到 \(\varprojlim_{n}\pi_{1}(\mathrm{P}_{n},0)\) 的由族 \(((r_n)_*)\) 得到的典范同态是单射，并且 \(\pi_1(\mathrm{P},0)\) 上的容许拓扑是 \(\varprojlim_{n}\pi_{1}(\mathrm{P}_{n},0)\) 上的逆极限拓扑的逆像；其中各群 \(\pi_1(\mathrm{P}_n,0)\) 都赋予离散拓扑。
\item
  证明 \(\Omega_{0}(\mathrm{P})\) 的道路连通分支是闭的，\(\Omega_{(0,0)}(\mathrm{P} \times \mathrm{P})\) 的道路连通分支也是闭的。
\end{enumerate}

\begin{enumerate}
\def\labelenumi{\arabic{enumi})}
\setcounter{enumi}{6}
\tightlist
\item
  我们沿用习题 6 的记号。
\end{enumerate}

\begin{enumerate}
\def\labelenumi{\alph{enumi})}
\item
  对每个 \(n \geqslant 1\)，令 \(\alpha_{n} \in \pi_{1}(P,0)\) 为 \(C_{n}\) 中一个圈的同伦类，且该同伦类生成 \(\pi_{1}(C_{n},0)\)。置 \(\beta_{n} = \alpha_{1}\alpha_{n}\alpha_{1}^{-1}\alpha_{n}^{-1}\)，\(\gamma_{n} = \beta_{n}^{n}\)。序列 \((\gamma_{n})\) 对容许拓扑收敛于 \(\varepsilon_{0}\)。
\item
  令 \(n \in N^{*}\)。对所有 \(c \in \Omega_{0}(P)\)，令 \(\omega_{n}(c)\) 为满足如下条件的整数 k 的上确界：存在由 I 中元素组成的严格递增序列 \((t_{1}, \ldots, t_{2k})\)，使得 \(c(t_{2i}) = 0\)，且 \(c(t_{2i-1}) = x_{n}\)，其中 \(i \in \{1, \ldots, k\}\)。证明 \(\omega_{n}(c)\) 是一个整数。证明映射 \(\omega_{n}\) 从 \(\Omega_{0}(P)\) 到 N 是上半连续的。
\item
  对每个 \(n \geqslant 1\)，令 \(c_{n}\) 为同伦类为 \(\gamma_{n}\) 的道路。证明 \(\omega_{1}(c_{n}) \geqslant 4n\)。推出序列 \((c_{n})\) 在 \(\mathcal{C}_{c}(\mathbf{I};\mathbf{P})\) 中没有极限。d) 容许拓扑下的群 \(\pi_{1}(\mathrm{P},0)\) 是豪斯多夫的、非离散的，并且严格粗于紧收敛拓扑。*空间 P 尤其不是可解圈的。*
\end{enumerate}

\begin{enumerate}
\def\labelenumi{\arabic{enumi})}
\setcounter{enumi}{7}
\tightlist
\item
  我们沿用习题 6 和 7 的记号；此外，令 \(p\colon \Omega_{0}(\mathrm{P})\to\pi_1(\mathrm{P},0)\) 为典范满射。
\end{enumerate}

\begin{enumerate}
\def\labelenumi{\alph{enumi})}
\item
  对每个整数对 \((m,n)\)，其中整数 \(\geqslant 1\) 且 \(m \neq n\)，置 \(\lambda_{m,n} = (\alpha_{m}\alpha_{n}\alpha_{m}^{-1}\alpha_{n}^{-1})^{m+n}\)，并置 \(\lambda_{m,n} = (\alpha_{1}\alpha_{m}\alpha_{1}^{-1}\alpha_{m}^{-1})^{n}\)；令 S 为形如 \((\lambda_{m,n},\mu_{m,n})\) 的 \(\pi_{1}(P,0)\times\pi_{1}(P,0)\) 中的二元组所成的集合。证明 S 不是闭集，但它在 \(\Omega_{0}(P)\times\Omega_{0}(P)\) 中的原像是闭的。推出映射 \(p\times p\colon\Omega_{0}(P)^{2}\to\pi_{1}(P,0)^{2}\) 不是严格的。
\item
  令 T 为形如 \(\lambda_{m,n}\mu_{m,n}\) 的 \(\pi_1(\mathrm{P},0)\) 中的圈类所成的集合。证明 T 是闭集。
\item
  证明 \(\pi_{1}(P,0)\) 的复合律在群 \(\pi_{1}(P,0)\) 赋予紧收敛拓扑的商拓扑时不是连续的。\(^{(5)}\)
\end{enumerate}

\begin{enumerate}
\def\labelenumi{\arabic{enumi})}
\setcounter{enumi}{8}
\tightlist
\item
  我们沿用习题 6 的记号。对每个整数 \(n \in \mathbf{N}^*\)，令 \(c_n: \mathbf{I} \to \mathrm{C}_n\) 为原点处由 \(t \mapsto \frac{1}{2n}(1 - \cos(2\pi t), \sin(2\pi t))\) 定义的圈。\\
\end{enumerate}

\begin{enumerate}
\def\labelenumi{\alph{enumi})}
\item
  对任意序列 \(\alpha \in \{0,1\}^{\mathbf{N}^*}\)，令 \(c_\alpha\) 为如下定义的映射：\(c_\alpha(0) = 0\)，并且 \(c_\alpha(t) = \alpha_n c_n (2^n t - 1)\)，其中 n 是满足 \(n \geqslant 1\) 的整数，且 \(t \in \mathbf{I}\) 满足 \(2^{-n} \leqslant t \leqslant 2^{1-n}\)。证明这是原点处的一个圈。证明映射 \(\alpha \mapsto [c_\alpha]\) 从 \(\{0,1\}^{\mathbf{N}^*}\) 到 \(\pi_1(\mathrm{P},0)\) 是单射。
\item
  对每个整数 \(n \geqslant 1\)，令 \(q_{n}\) 为从 P 到 \(C_{n}\) 的映射：有 \(q_{n}(x) = x\)（\(x \in C_{n}\)），否则 \(q_{n}(x) = 0\)；它是连续的。证明群同态 \(\psi: \pi_{1}(\mathrm{P}, 0) \to \prod_{n \geqslant 1} \pi_{1}(\mathrm{C}_{n}, 0)\)，由 \(\psi(\gamma) = ((q_{n})_{*}(\gamma))_{n}\) 给出，具有一个截面。
\end{enumerate}

\begin{enumerate}
\def\labelenumi{\arabic{enumi})}
\setcounter{enumi}{9}
\tightlist
\item
  对每个整数 \(n \geqslant 1\)，令 \(\mathrm{C}_n\) 为数值平面中以 \((1/n,0)\) 为圆心、半径为 \(1/n\) 的圆；令 P 为所有空间 \(\mathrm{C}_n\)（\(n \geqslant 1\)）的并。对每个整数 \(n \geqslant 1\)，令 \(\mathrm{B}_n\) 为平面上的点 \((x,y)\) 所成的集合，其中满足 \((nx-1)^2 + n^2y^2 \leqslant 1 \leqslant ((n+1)x-1)^2 + (n+1)^2y^2\)；令 B 为平面上的点 \((x,y)\) 所成的集合，其中满足 \((x-1)^2 + y^2 \leqslant 1\)。置 \(\mathrm{B}_0 = \mathrm{P}\)，令 A 为将族 \((\mathrm{B}_n)_{n \geqslant 0}\) 的和空间按如下最粗等价关系取商所得的空间：识别 \((b,n)\) 与 \((b,0)\)，其中 \(b \in \mathrm{C}_n \cup \mathrm{C}_{n+1}\)。
\end{enumerate}

\begin{enumerate}
\def\labelenumi{\alph{enumi})}
\item
  证明存在唯一的连续映射 \(p \colon \mathrm{A} \to \mathrm{B}\)，对每个整数 \(n \geqslant 0\) 及每个 \(b \in \mathrm{B}_n\)，该映射都把 \((b, n)\) 的等价类映到点 \(b\)。证明 \(p\) 是双射但不是同胚。
\item
  证明空间 A 是局部道路连通的。
\item
  令 a 为点 \((0,0)\) 的等价类。证明 \(\pi_{1}(A,a)\) 上的紧收敛拓扑的商拓扑是粗拓扑。
\item
  证明空间 A 是单连通的。
\item
  令 j 为从 \(C = B_{0}\) 到 A 的典范映射。证明映射 \(j_{*} \colon \pi_{1}(C,0) \to \pi_{1}(A,j(0))\) 是满射，且其核的基数为连续统。
\item
  证明群 \(\pi_{1}(A,j(0))\) 的基数为连续统。*特别地，空间 A 不是可解圈的。*
\end{enumerate}

